\section{Conclusiones}

\subsection{Conclusión General}

El desarrollo de la base de datos ECOBICI permitió aplicar de manera integral los conocimientos adquiridos durante el curso de Bases de Datos, abarcando las etapas de análisis, diseño conceptual, diseño lógico e implementación física. A lo largo del proyecto se construyó una solución capaz de administrar la información relacionada con usuarios, membresías, estaciones, bicicletas, viajes, incidentes y personal operativo, garantizando la integridad y consistencia de los datos mediante restricciones, llaves primarias, llaves foráneas y mecanismos de seguridad.

Asimismo, la implementación de procedimientos almacenados, funciones, vistas, triggers e informes estadísticos permitió automatizar procesos y facilitar la obtención de información relevante para la toma de decisiones. El proyecto permitió comprender la importancia de un diseño adecuado desde las primeras etapas de desarrollo, demostrando que una base de datos correctamente normalizada contribuye significativamente a la eficiencia, mantenimiento y escalabilidad de los sistemas de información.

\subsection{Conclusión de David Díaz Antúnez}

La realización de este proyecto permitió reforzar los conocimientos relacionados con el diseño e implementación de bases de datos relacionales, así como comprender la importancia de una correcta organización de la información para garantizar la integridad de los datos. Durante el desarrollo del sistema ECOBICI se adquirió experiencia en el uso de SQL Server, la creación de estructuras relacionales y la aplicación de buenas prácticas de modelado. Además, se fortalecieron habilidades de trabajo colaborativo y resolución de problemas, aspectos fundamentales en el desarrollo de proyectos de software.

\subsection{Conclusión de Mauricio Gabriel Hernández Acosta}

A través del desarrollo de la base de datos ECOBICI fue posible aplicar de manera práctica los conceptos estudiados durante el curso, especialmente aquellos relacionados con la normalización, la integridad referencial y la implementación de objetos programados. El proyecto permitió comprender cómo las decisiones tomadas durante la etapa de diseño impactan directamente en el rendimiento y mantenimiento de una base de datos. Asimismo, se obtuvo experiencia en la construcción de consultas complejas e informes estadísticos orientados al análisis de información y apoyo en la toma de decisiones.

\subsection{Conclusión de César Ricardo Sánchez Luján}

El proyecto ECOBICI representó una oportunidad para integrar conocimientos teóricos y prácticos en un entorno cercano a una situación real. La implementación de mecanismos de seguridad, restricciones de integridad y procesos automatizados permitió comprender la importancia de diseñar soluciones robustas y confiables. De igual forma, se fortaleció la capacidad para analizar requerimientos, modelar información y desarrollar bases de datos que satisfagan necesidades específicas de una organización, contribuyendo a la formación profesional en el área de tecnologías de la información.
